\documentclass[12pt]{article}
\usepackage{graphicx} % Required for inserting images
\usepackage[margin=0.5in]{geometry}
\usepackage{amsmath, amssymb}
\usepackage{mathrsfs}


\usepackage{soul}


% Dr. David Brown Commands
\newcommand{\ds}{\displaystyle}
\newcommand{\R}{\mathbb{R}}
\newcommand{\N}{\mathbb{N}}
\newcommand{\Q}{\mathbb{Q}}
\newcommand{\C}{\mathbb{C}}


% Commands to get script r
\usepackage{calligra}
\DeclareMathAlphabet{\mathcalligra}{T1}{calligra}{m}{n}
\DeclareFontShape{T1}{calligra}{m}{n}{<->s*[2.2]callig15}{}
\newcommand{\scripty}[1]{\ensuremath{\mathcalligra{#1}}}

% Unit commands
\newcommand{\degree}{\text{°}}
\newcommand{\unitSpeed}{\frac{\text{m}}{\text{s}}}
\newcommand{\unitAcceleration}{\frac{\text{m}}{\text{s}^2}}

% Constant commands
\newcommand{\avogadro}{6.022\times10^{23}}

\newcommand{\boltzmannConstK}{1.381\times10^{-23}\frac{\text{J}}{\text{K}}}
\newcommand{\boltzmannConstInEV}{8.617\times10^{-5}\frac{\text{eV}}{\text{K}}}
\newcommand{\planckConst}{6.626\times10^{-34}\text{J}\cdot\text{s}}

\newcommand{\atmP}{1.01325\times10^5\,\text{Pa}}

% Known Value Commands
\newcommand{\electronMass}{9.109\times10^{-31}\,\text{kg}}
\newcommand{\protonMass}{1.673\times10^{-27}\,\text{kg}}

% Commands to easily write partial derivatives
\newcommand{\partDeriv}[2]{\frac{\partial #1}{\partial #2}}
\newcommand{\doublePartDeriv}[2]{\frac{\partial^2 #1}{\partial^2 #2}}


\title{Problem 7.39}
\author{Gabriel Decker}


\begin{document}


\maketitle
\begin{flushleft}


In this problem, we consider the Planck Spectrum.
We will begin with the equation of the total photon energy per unit area as a function of photon energy $\epsilon$ and do a change of variables with $\lambda=hc/\epsilon$.
This will change the integrand's original interpretation of energy density per photon energy to energy density per photon wavelength.\\~\\

We begin with the given equation that integrates over photon energy.
\begin{equation}
    \frac{U}{V}=\int_0^\infty\frac{8\pi}{(hc)^3}\frac{\epsilon^3}{e^{\epsilon/kT}-1}\,d\epsilon
\end{equation}
With $\lambda=hc/\epsilon$, this implies that $\epsilon=hc/\lambda$.
This implies that $d\epsilon/d\lambda=-hc/\lambda^2$, so $d\epsilon=(-hc/\lambda^2)/d\lambda$.
We can now apply this change of variables to the equation.
Let's begin by just substituting $\epsilon$.
$$\frac{U}{V}=\frac{8\pi}{(hc)^3}\int_0^\infty\frac{\epsilon^3}{e^{\epsilon/kT}-1}\,d\epsilon$$
$$\implies\frac{U}{V}=\frac{8\pi}{(hc)^3}\int_0^\infty\frac{(hc/\lambda)^3}{e^{(hc/\lambda)/kT}-1}\,d\epsilon$$
$$\implies\frac{U}{V}=8\pi\int_0^\infty\frac{1}{\lambda^3}\frac{1}{e^{hc/\lambda kT}-1}\,d\epsilon$$
Replacing the variable of integration, we get the following.
$$\implies\frac{U}{V}=8\pi\int_\infty^0\frac{1}{\lambda^3}\frac{1}{e^{hc/\lambda kT}-1}\cdot\left(-\frac{hc}{\lambda^2}\right)\,d\lambda$$
$$\implies\frac{U}{V}=8\pi hc\int_0^\infty\frac{1}{\lambda^5}\frac{1}{e^{hc/\lambda kT}-1}\,d\lambda$$
This makes the photon wavelength spectrum the following.
\begin{equation}
    u(\lambda)=\frac{8\pi hc}{\lambda^5}\frac{1}{e^{hc/\lambda kT}-1}
\end{equation}\\~\\

Now that we have an expression for the energy density per photon wavelength distribution, we can learn some things about it.
We can plot the distribution at, for example $T=300\,\text{K}$ and find where the max value is.
I'll use sympy, numpy, and matplotlib to accomplish this.
I used a substitution of a dimensionless variable: $x=\lambda kT/hc$.
So, $\lambda=xhc/kT$.\\
TODO: Getting not implemented error. Try defining $x$ such that we get $e^x$ instead of $e^{1/x}$, and then invert the solution instead or something like that.




\end{flushleft}

\end{document}
